\section{Small-arm inputs and the short four-arm theorem}\label{shared:short-cores}
\subsection{Construction inputs for small arms}
\label{seven:small-inputs}

The following are the precise small-arm inputs needed in this section.
Their proofs and their complete finite tables are in
Sections~\ref{app:small-alpha}, \ref{app:pair-offsets}, and
\ref{app:finite-bases}.

\begin{proposition}\label{seven:small-alpha}
The spider $\Sp{s,a,b}$ has a root-zero alpha labeling in each of the
following cases:
\begin{enumerate}
\item $s=1$, provided $(a,b)\ne(2,2)$;
\item $s\in\{2,3,4,5\}$, $s\le a\le b$, and $b\ge7$;
\item $s\in\{2,3,4,5\}$ and $a=b=6$.
\end{enumerate}
\end{proposition}

The first assertion follows by adjoining a new maximum leaf to the
prescribed-zero alpha path on the other two arms.  For the second,
the appendix supplies explicit near-band constructions and their
resonance repairs.  To specify exactly what those proofs must establish,
put $e=\floor{s/2}$ and $k=e+\floor{a/2}$.  The alpha extension lemma
already gives every $b>2k$ other than $b=4k+1$.  Thus the additional
construction obligations are exactly
\begin{center}
\begin{tabular}{c|c|c|c}
$s$ & parity of $a$ & near band & resonance\\ \hline
$2,3$ & odd & $a\le b\le a+1$ & $b=2a+3$\\
$2,3$ & even & $a\le b\le a+2$ & $b=2a+5$\\
$4,5$ & odd & $a\le b\le a+3$ & $b=2a+7$\\
$4,5$ & even & $a\le b\le a+4$ & $b=2a+9$.
\end{tabular}
\end{center}
Only values with $b\ge7$ are asserted in the second part.  The
small-arm appendices prove these bands by explicit path formulas,
fixed affine edge trades with all-parameter topology checks, and their
exhaustive finite complements.  In particular, a finite experiment
inside a near band is not used as an existence argument outside it.
The third part consists of the four stated literal alpha labelings.

\begin{proposition}\label{seven:small-pairs}
For $x\in\{2,3,4,5\}$ and $n\ge x$,
\[
   \{0,\ldots,x-1\}\subseteq\Om(\Sp{x,n}).
\]
In addition, $2\in\Om(\Sp{2,n})$ for every $n\ge3$.
\end{proposition}

The pair appendix gives the cracked-permutation formulas, the ordinary
graceful tails, and every finite starting word needed here.  The prescribed-first ordinary path theorem supplies each graceful
tail in the stated range; the displayed prefixes partition all labels and colors.  Consequently these are uniform
statements for every $n$ in the displayed ranges.

We also use the hub-zero three-arm theorem
(Theorem~\ref{common:hubzero-three}).  The finite basis denoted below by
$F_4$ consists of all $\binom84=70$ sorted four-tuples with
entries in $\{2,\ldots,6\}$, each accompanied by a hub-zero graceful
labeling.  Verification of a finite labeling is explained in
Section~\ref{templates:finite-check}; the full words belong to
Section~\ref{app:finite-bases}.

\subsection{Four arms with a short arm}
\label{seven:four-short-section}

\begin{lemma}\label{seven:four-short}
Every four-arm spider with an arm of length at most five has a
hub-zero graceful labeling.
\end{lemma}
\begin{proof}
An arm of length one can be added by a new maximum leaf to the
hub-zero labeling of the other three arms.  It therefore remains to
treat minimum lengths two, three, four, and five.  In each case, if
the largest length is at most six, $F_4$ applies.  Hence throughout
the remaining argument the largest length $c$ is at least seven.

\medskip\noindent
\emph{Minimum two.}
Write the lengths as $2\le a\le b\le c$, with an additional arm of
length two.  Proposition~\ref{seven:small-alpha}, together with the
nine literal small positive cores in the finite appendix, supplies
$\Sp{2,a,b}$ except possibly when
\begin{equation}\label{seven:min2-exceptions}
 (a,b)\in\{(2,2),(2,3),(2,4),(2,6),(3,3),(4,4)\}.
\end{equation}
The nine positive cores are
\[
 (2,2,5),\ (2,3,4),\ (2,3,5),\ (2,3,6),\ (2,4,5),\
 (2,4,6),\ (2,5,5),\ (2,5,6),\ (2,6,6).
\]
No nonexistence claim about the exceptional cores is required.

For a supplied alpha core its index is
$J=1+\floor{a/2}+\floor{b/2}\le b+1$.  If $c\ge b+2$, then
$c>J$ and ordinary attachment applies.  If $c=b$, necessarily $b\ge7$.
When $a=2$, $J=2+\floor{b/2}<b$, so attachment still applies.
When $a\ge3$, use instead the alpha pair $\Sp{2,a}$, whose index
$1+\floor{a/2}\le b-1$ lies in the equal-pair offset interval for
$\Sp{b,b}$.  If $c=b+1$, attachment applies except when $a=b$ is
even.  In that remaining case $b\ge6$; the alpha pair $\Sp{2,c}$
has index $1+b/2\le b-1$ and can be combined with the equal pair
$\Sp{b,b}$.  This also includes $b=6,c=7$.

For the pairs $(2,2),(2,4),(2,6),(4,4)$ in
\eqref{seven:min2-exceptions}, at most one of the four arms is odd.
Halving every even length and flooring the possible odd length gives
a four-arm contree with a unit arm.  That contree is hub-zero graceful
by the first paragraph.  The ordinary full-subdivision or
perfect-matching lift of the foundations gives the required
root-zero labeling of the target.  For $(a,b)=(2,3)$, compose the
alpha pair $\Sp{2,3}$, of index two, with the offset-two pair
$\Sp{2,c}$.  For $(a,b)=(3,3)$, use the same alpha pair with the
offset-two pair $\Sp{3,c}$.  Proposition~\ref{seven:small-pairs}
provides both offsets.  This exhausts minimum two.

\medskip\noindent
\emph{Minimum three.}
Write the spider as $\Sp{3,a,b,c}$ with $3\le a\le b\le c$.
The supplied positive alpha cores $\Sp{3,a,b}$ include every case
except possibly $(a,b)=(3,3)$ and $(4,4)$.  For $b\le6$, the
eight positive literal cores are
\[
 (3,3,4),\ (3,3,5),\ (3,3,6),\ (3,4,5),\
 (3,4,6),\ (3,5,5),\ (3,5,6),\ (3,6,6).
\]
For a positive core, $J=1+\floor{a/2}+\floor{b/2}\le b+1$.
Thus $c\ge b+2$ is an ordinary attachment.  If $c=b$, the index
$1+\floor{a/2}$ of $\Sp{3,a}$ is at most $b-1$; combine it with
the equal offset pair $\Sp{b,b}$.  If $c=b+1$, attachment fails
only when $a=b$ is even.  Then $b\ge6$ and the alpha pair
$\Sp{3,c}$ has index $1+b/2\le b-1$, so the same equal-pair
composition applies.  For the exceptional core $(a,b)=(3,3)$, use
the alpha pair $\Sp{3,3}$ of index two and the offset-two pair
$\Sp{3,c}$.  For $(a,b)=(4,4)$, use the alpha pair $\Sp{3,4}$ of
index three and the offset-three pair $\Sp{4,c}$.  These offsets
are covered by Proposition~\ref{seven:small-pairs}.

\medskip\noindent
\emph{Minimum four or five.}
Write $\Sp{s,a,b,c}$, where $s\in\{4,5\}$ and
$s\le a\le b\le c$.  The small-alpha theorem and the six literal
cores
\[
 (4,5,5),\ (4,5,6),\ (4,6,6),\
 (5,5,5),\ (5,5,6),\ (5,6,6)
\]
give an alpha core $\Sp{s,a,b}$ except possibly for
\begin{equation}\label{seven:min4-exceptions}
 (s,a,b)=(4,4,4),\ (4,4,5),\ (4,4,6).
\end{equation}
For the first and third exceptional cores, there is at most one odd
arm in the target four-arm spider.  Its halved contree contains a
two-arm, so the already established minimum-two case and the ordinary
parity lift apply.  The same argument covers $\Sp{4,4,5,c}$ when
$c$ is even.  When $c$ is odd, compose the alpha pair $\Sp{4,4}$
of index four with the offset-four pair $\Sp{5,c}$; here $c\ge7$,
and Proposition~\ref{seven:small-pairs} applies.

Now suppose that the alpha core is supplied.  Put
\[
 k=2+\floor{a/2},\qquad J=k+\floor{b/2}.
\]
If $c>J$, attach $c$.  Hence assume $c\le J$.  We have $k\le b-1$:
the only possible failure with $s\le a\le b$ would be $a=b=4$,
already in \eqref{seven:min4-exceptions}.  If $c=b$, use the alpha
pair $\Sp{s,a}$ and the equal offset pair $\Sp{b,b}$.
Assume therefore $c>b$.  The inequalities
$b\le c\le k+\floor{b/2}$ give $\ceil{b/2}\le k$, and hence
$\floor{c/2}\le k$.  The HIGH pair lemma then gives the desired
offset $k$ on $\Sp{b,c}$ unless
\begin{equation}\label{seven:small-resonance}
 c=4q+1,\qquad k=b-1-q
\end{equation}
for an integer $q\ge1$.

We reduce this resonance exactly.  Equations
\eqref{seven:small-resonance} and $c\le J$ imply
\[
 \floor{a/2}=b-q-3,\qquad
 5q+2\le b+\floor{b/2}.
\]
Since $a\le b$, $b\le2q+6$.  Therefore
\[
 5q+2\le b+\floor{b/2}\le3q+9,
\]
which forces $q\le3$.  If $q=1$, then $c=5>b$ forces $b\le4$,
whereas $7\le b+\floor{b/2}$ forces $b\ge5$, a contradiction.
If $q=2$, then $c=9$, $b=8$, and $a\in\{6,7\}$.
If $q=3$, then $c=13$, $b=12$, and $a=12$.

The repeated case $a=b$ has another equal-pair composition.  In this
case $q=b-3-\floor{b/2}$, and the alpha pair $\Sp{s,c}$ has index
$2+2q$, equal to $b-4$ for even $b$ and $b-3$ for odd $b$.
This is an available offset on $\Sp{b,b}$.  Consequently only
\[
 \Sp{s,6,8,9},\qquad \Sp{s,7,8,9},\qquad s\in\{4,5\},
\]
remain.  For $\Sp{s,6,8,9}$ use the alpha pair $\Sp{s,8}$ of
index six and the following offset-six pair on $\Sp{6,9}$:
\begin{align*}
 &(15,13,14,20,7,21),\\
 &(12,16,11,8,19,9,18,10,17).
\end{align*}
Its labels are $[7,21]$ and its colors are $[1,15]$.  It can also
be obtained by factor-three offset subdivision of the offset-two
pair $(5,7)\mid(4,3,6)$ on $\Sp{2,3}$.

For $\Sp{4,7,8,9}$, the following words, read from hub zero, have
lengths $4,9,8,7$:
\begin{align*}
 &(28,1,27,2),\\
 &(24,3,26,4,23,5,25,8,22),\\
 &(16,6,21,9,20,11,19,12),\\
 &(13,7,10,15,17,18,14).
\end{align*}
They use every label and color from one through $28$ once.
Replace the first word by $(29,1,28,2,27)$ to obtain
$\Sp{5,7,8,9}$: the short color interval $[25,28]$ becomes
$[25,29]$, and all other edges are unchanged.  This finishes the
last two profiles, and the proof of the lemma.
\end{proof}

